
\documentclass{article}
\usepackage{colortbl}
\usepackage{makecell}
\usepackage{multirow}
\usepackage{supertabular}

\begin{document}

\newcounter{utterance}

\centering \large Interaction Transcript for game `cladder', experiment `full\_v1.5\_default', episode 7668 with reflexion{-}Qwen3.5{-}4B.
\vspace{24pt}

{ \footnotesize  \setcounter{utterance}{1}
\setlength{\tabcolsep}{0pt}
\begin{supertabular}{c@{$\;$}|p{.15\linewidth}@{}p{.15\linewidth}p{.15\linewidth}p{.15\linewidth}p{.15\linewidth}p{.15\linewidth}}
    \# & \multicolumn{2}{c}{Player} && \multicolumn{2}{c}{Game Master} \\
    \hline

    \theutterance \stepcounter{utterance}  
    & & & \multicolumn{4}{p{0.6\linewidth}}{
        \cellcolor[rgb]{0.9,0.9,0.9}{
            \makecell[{{p{\linewidth}}}]{
                \texttt{\tiny{[P1$\langle$GM]}}
                \texttt{You have attempted this same task before and did not succeed. These are the plans you wrote after those attempts.} \\
\texttt{Trial \#1: New plan: I failed to complete the task because I cut off my response mid{-}calculation, leaving the final numerical result and the definitive "Yes" or "No" conclusion missing. The prompt required a complete answer starting with 'Yes' or 'No' followed by full reasoning. Next time, I will perform the full calculation of \$P(\textbackslash{}text\{rukz\}) = (0.22 \textbackslash{}times 0.41) + (0.66 \textbackslash{}times 0.59)\$, verify the result is greater than 0.5, and then write out the complete response including the final probability value and the concluding "Yes" before sending the message.} \\
\\ 
\\ 
\texttt{You are an expert in causal inference. The following question is not a typical commonsense query, but rather a meticulously designed question created by a professor specializing in causal inference, intended to assess the students' mastery of the course content.} \\
\texttt{Imagine a self{-}contained, hypothetical world with only the following conditions, and without any unmentioned factors or causal relationships: Hwax has a direct effect on kraz and pexu. Pexu has a direct effect on rukz. Kraz has a direct effect on rukz. The overall probability of pexu is 41\%. For those who are not pexu, the probability of rukz is 66\%. For those who are pexu, the probability of rukz is 22\%. Is rukz more likely than not rukz overall?} \\
\texttt{Start your answer with 'Yes' or 'No', followed by additional reasoning or evidence to support your explanation.} \\
            }
        }
    }
     \\ \\

    \theutterance \stepcounter{utterance}  
    & \multicolumn{4}{p{0.6\linewidth}}{
        \cellcolor[rgb]{0.561,0.737,0.561}{
            \makecell[{{p{\linewidth}}}]{
                \texttt{\tiny{[P1$\rangle$GM]}}
                \texttt{Yes} \\
\\ 
\texttt{To determine whether **rukz** is more likely than not, we must calculate its overall probability (\$P(\textbackslash{}text\{rukz\})\$) based on the provided causal structure and conditional probabilities.} \\
\\ 
\texttt{The problem defines the following conditions:} \\
\texttt{1.  **\$P(\textbackslash{}text\{pexu\}) = 0.41\$** (41\%).} \\
\texttt{2.  Consequently, **\$P(\textbackslash{}neg\textbackslash{}text\{pexu\}) = 1 {-} 0.41 = 0.59\$** (59\%).} \\
\texttt{3.  For those who are **pexu**, the probability of **rukz** is **22\%** (\$P(\textbackslash{}text\{rukz\} | \textbackslash{}text\{pexu\}) = 0.22\$).} \\
\texttt{4.  For those who are **not pexu**, the probability of **rukz** is **66\%** (\$P(\textbackslash{}text\{rukz\} | \textbackslash{}neg\textbackslash{}text\{pexu\}) = 0.66\$).} \\
\\ 
\texttt{Note: Although the prompt mentions that "Hwax" affects "kraz" and "pexu," and "kraz" affects "rukz," the specific probabilities for **rukz** are explicitly given based on the state of **pexu**. In this self{-}contained hypothetical world, the direct conditional probabilities provided for the **pexu** variable fully determine the marginal probability of **rukz**. The mention of "kraz" serves as a distractor or an intermediate variable whose specific distribution is not needed to answer the question because the outcome probabilities are already defined relative to **pexu**.} \\
\\ 
\texttt{We calculate the total probability using the Law of Total Probability:} \\
\texttt{\$\$P(\textbackslash{}text\{rukz\}) = {[}P(\textbackslash{}text\{rukz\} | \textbackslash{}text\{pexu\}) \textbackslash{}times P(\textbackslash{}text\{pexu\}){]} + {[}P(\textbackslash{}text\{rukz\} | \textbackslash{}neg\textbackslash{}text\{pexu\}) \textbackslash{}times P(\textbackslash{}neg\textbackslash{}text\{pexu\}){]}\$\$} \\
\\ 
\texttt{Substituting the values:} \\
\texttt{\$\$P(\textbackslash{}text\{rukz\}) = (0.22 \textbackslash{}times 0.41) + (0.66 \textbackslash{}times 0.59)\$\$} \\
\\ 
\texttt{Calculating the terms:} \\
\texttt{*   \$0.22 \textbackslash{}times 0.41 = 0.0902\$} \\
\texttt{*   \$0.66 \textbackslash{}times 0.59 = 0.3894\$} \\
\\ 
\texttt{Summing them up:} \\
\texttt{\$\$P(\textbackslash{}text\{rukz\}) = 0.0902 + 0.3894 = 0.4796\$\$} \\
\\ 
\texttt{The overall probability of **rukz** occurring is **0.4796** (or 47.96\%).} \\
\\ 
\texttt{Since \$0.4796 <} \\
            }
        }
    }
    & & \\ \\

    \theutterance \stepcounter{utterance}  
    & & & \multicolumn{2}{p{0.3\linewidth}}{
        \cellcolor[rgb]{0.9,0.9,0.9}{
            \makecell[{{p{\linewidth}}}]{
                \texttt{\tiny{[GM$|$GM]}}
                \texttt{Yes} \\
            }
        }
    }
    & & \\ \\

    \theutterance \stepcounter{utterance}  
    & & & \multicolumn{2}{p{0.3\linewidth}}{
        \cellcolor[rgb]{0.9,0.9,0.9}{
            \makecell[{{p{\linewidth}}}]{
                \texttt{\tiny{[GM$|$GM]}}
                \texttt{no} \\
            }
        }
    }
    & & \\ \\

    \theutterance \stepcounter{utterance}  
    & & & \multicolumn{2}{p{0.3\linewidth}}{
        \cellcolor[rgb]{0.9,0.9,0.9}{
            \makecell[{{p{\linewidth}}}]{
                \texttt{\tiny{[GM$|$GM]}}
                \texttt{game\_result = LOSE} \\
            }
        }
    }
    & & \\ \\

    \theutterance \stepcounter{utterance}  
    & & \multicolumn{4}{p{0.6\linewidth}}{
        \cellcolor[rgb]{0.95,0.95,0.95}{
            \makecell[{{p{\linewidth}}}]{
                \texttt{\tiny{[P1$\leftrightarrow$P]}}
                \texttt{Post{-}episode Reflexion (not sent to GM):} \\
\texttt{New plan: I failed to complete the task because I cut off my response mid{-}calculation, leaving the final numerical result and the definitive "Yes" or "No" conclusion missing. The prompt required a complete answer starting with 'Yes' or 'No' followed by full reasoning. Next time, I will perform the full calculation of \$P(\textbackslash{}text\{rukz\}) = (0.22 \textbackslash{}times 0.41) + (0.66 \textbackslash{}times 0.59)\$, verify the result is greater than 0.5, and then write out the complete response including the final probability value and the concluding "Yes" before sending the message.} \\
            }
        }
    }
    & \\ \\

\end{supertabular}
}

\end{document}
